% Plantilla mínima de los boletines IES (pandoc --template).
% Solo carga lo estrictamente necesario por el contenido; la tipografía y la
% geometría las decide LaTeX (tamaño carta, 12pt, márgenes por defecto).

\documentclass[12pt,letterpaper]{article}

% Codificación y acentos (pdflatex)
\usepackage[T1]{fontenc}
\usepackage[utf8]{inputenc}

% Fuente escalable (Latin Modern); evita las bitmap EC de METAFONT
\usepackage{lmodern}

% Español: títulos y pies ("Figura"), guionado
\usepackage[spanish]{babel}

% Secciones sin numeración automática (equivalente a \section*{...}): solo se
% muestran los números que el markdown escribe a mano ("1.", "2.", "A.", ...).
\setcounter{secnumdepth}{-1}

% Los títulos nunca quedan huérfanos al pie de página: reserva espacio antes
% de cada sección y, si no cabe el título + unas líneas, salta a la página
% siguiente. Necesario sobre todo cuando el título precede a una tabla longtable.
\usepackage{needspace}
\usepackage{etoolbox}
\pretocmd{\section}{\needspace{3\baselineskip}}{}{}
\pretocmd{\subsection}{\needspace{3\baselineskip}}{}{}
\pretocmd{\subsubsection}{\needspace{3\baselineskip}}{}{}

% Párrafos: penalizar líneas huérfanas (club) y viudas (widow) al cambiar de
% página: la primera línea de un párrafo no queda sola abajo, ni la última sola arriba.
\clubpenalty=10000
\widowpenalty=10000
\displaywidowpenalty=10000

% Matemáticas (variables como T, I, neq)
\usepackage{amsmath,amssymb}

% Imágenes: los diagramas PDF se ajustan al ancho de caja, sin anchos fijos
\usepackage{graphicx}
\makeatletter
\def\maxwidth{\ifdim\Gin@nat@width>\linewidth\linewidth\else\Gin@nat@width\fi}
\def\maxheight{\ifdim\Gin@nat@height>\textheight\textheight\else\Gin@nat@height\fi}
\makeatother
\setkeys{Gin}{width=\maxwidth,height=\maxheight,keepaspectratio}

% Fondo: membrete a página completa (carta) en TODAS las hojas. La ruta la
% inyecta build-boletin.sh con -V membrete=<ruta absoluta>; si la variable no
% se define, el fondo simplemente no se carga. Ojo: en un template de pandoc
% los signos de dolar se interpretan como variables de template, incluso en los
% comentarios; por eso aqui no debe aparecer ningun signo de dolar suelto
% aparte de la variable membrete del bloque condicional.
\usepackage{background}
% NOTA: hay que dar width Y height explicitos (paperwidth x paperheight). Las
% claves globales de Gin (width=\maxwidth, height=\maxheight, keepaspectratio)
% fijadas mas arriba para los diagramas se aplican por defecto a todo
% \includegraphics; con solo width=\paperwidth el fondo quedaria reducido al
% area de texto (margenes). Con ambas dimensiones explicitas y la misma
% proporcion carta del membrete, llena la pagina sin deformar ni recortar.
\backgroundsetup{
  scale=1,
  angle=0,
  opacity=1,
  contents={\includegraphics[width=\paperwidth,height=\paperheight]{/home/kritt/Git/boletin/.opencode/skills/boletin-latex-pdf/assets/membrete.pdf}}
}

% Tablas (los boletines pueden traer tablas Markdown)
\usepackage{longtable,booktabs,array,calc}

% Pies de figura
\usepackage{caption}
\captionsetup{font=small,labelfont=bf}

% Listas compactas de Markdown
\providecommand{\tightlist}{%
  \setlength{\itemsep}{0pt}\setlength{\parskip}{0pt}}

% Enlaces y destinos (al final, como recomienda hyperref)
\usepackage{hyperref}

\begin{document}

\hypertarget{termodinuxe1mica-financiera-08}{%
\section{TERMODINÁMICA FINANCIERA
08}\label{termodinuxe1mica-financiera-08}}

\hypertarget{titular-y-resumen-ejecutivo-la-ilusiuxf3n-del-presupuesto-frente-a-la-restricciuxf3n-sistuxe9mica}{%
\subsection{1. Titular y Resumen Ejecutivo: La Ilusión del Presupuesto
frente a la Restricción
Sistémica}\label{titular-y-resumen-ejecutivo-la-ilusiuxf3n-del-presupuesto-frente-a-la-restricciuxf3n-sistuxe9mica}}

\href{https://www.eleconomista.com.mx/finanzaspersonales/educacion-financiera-funciona-dice-foro-economico-mundial-20261007-837387.html}{Educacion
financiera}

El Foro Económico Mundial (WEF) ha puesto sobre la mesa un debate
incómodo pero necesario: \textbf{la educación financiera tradicional no
funciona}. Pretender resolver la precariedad económica y el
endeudamiento masivo de los hogares mediante talleres de ahorro,
presupuestos y recortes individuales equivale a intentar mantener
cientos de pelotas de ping pong bajo el agua en una piscina presurizada.
El fallo no reside en una supuesta ignorancia del consumidor, sino en
una falla sistémica: el costo operativo vital (\(OE\)) supera con creces
el Throughput (\(T\)) salarial real disponible para la base de la
pirámide económica global.

\begin{center}\rule{0.5\linewidth}{0.5pt}\end{center}

\hypertarget{diagnuxf3stico-del-sistema-seguxfan-toc}{%
\subsection{2. Diagnóstico del Sistema según
TOC}\label{diagnuxf3stico-del-sistema-seguxfan-toc}}

Bajo la óptica de la \textbf{Teoría de Restricciones (TOC)}, el sistema
socioeconómico actual impone un cuello de botella físico y estructural
en la generación de ingresos reales para los individuos, mientras que
los costos de supervivencia (vivienda, salud, alimentos, energía)
escalan de manera inelástica.

\begin{figure}
\centering
\includegraphics{Diagramas/figura-1.pdf}
\caption{El Falso Dilema de la Educación Financiera y el Flujo
Deficitario}
\end{figure}

\begin{itemize}
\tightlist
\item
  \textbf{Throughput (\(T\)):} La velocidad de creación y captura de
  riqueza por parte del trabajador está limitada por estructuras de
  mercado asimétricas y salarios reales estancados.
\item
  \textbf{Inventario / Inversión (\(I\)):} Pasivos obligatorios (deudas
  por consumo básico y compromisos habitacionales) que actúan como
  anclas financieras inevitables.
\item
  \textbf{Gasto de Operación (\(OE\)):} El costo mínimo de existencia en
  el sistema, impulsado por dinámicas inflacionarias globales que
  escapan por completo al control del individuo.
\end{itemize}

\begin{center}\rule{0.5\linewidth}{0.5pt}\end{center}

\hypertarget{dicotomuxeda-clave-1-contabilidad-de-costos-tradicional-vs.-contabilidad-del-throughput-ta}{%
\subsection{3. Dicotomía Clave 1: Contabilidad de Costos Tradicional
vs.~Contabilidad del Throughput
(TA)}\label{dicotomuxeda-clave-1-contabilidad-de-costos-tradicional-vs.-contabilidad-del-throughput-ta}}

\begin{itemize}
\tightlist
\item
  \textbf{La Visión de la Contabilidad de Costos Tradicional:} Evalúa al
  individuo bajo la premisa errónea de la \emph{eficiencia local}. Lo
  trata como un centro de costos ineficiente al que hay que auditar,
  recortando gastos hormiga y minimizando partidas individuales (el
  café, las suscripciones, el ocio). Asume que la salud financiera es
  una simple ecuación aritmética de disciplina microeconómica, culpando
  directamente a la víctima del sistema por su falta de ``austeridad''.
\item
  \textbf{La Visión de la Contabilidad del Throughput (TA):} Comprende
  que el \(OE\) (gasto operativo) tiene un límite biológico inferior
  irreductible: no se puede recortar el costo de subsistencia por debajo
  de cero sin colapsar al agente. El foco de la TA dictamina que
  \textbf{la riqueza y la viabilidad no provienen de exprimir el \(OE\)
  localmente, sino de maximizar el \(T\) y eliminar las restricciones
  sistémicas} que limitan la capacidad de generar valor real en el
  mercado. Si el entorno estrangula el flujo de ingresos, cualquier
  manual de finanzas personales se vuelve operativamente estéril.
\end{itemize}

\begin{center}\rule{0.5\linewidth}{0.5pt}\end{center}

\hypertarget{dicotomuxeda-clave-2-teoruxeda-monetaria-moderna-mmt-vs.-teoruxeda-de-restricciones-toc}{%
\subsection{4. Dicotomía Clave 2: Teoría Monetaria Moderna (MMT)
vs.~Teoría de Restricciones
(TOC)}\label{dicotomuxeda-clave-2-teoruxeda-monetaria-moderna-mmt-vs.-teoruxeda-de-restricciones-toc}}

\begin{itemize}
\tightlist
\item
  \textbf{Visión MMT:} Analiza el problema desde la macroeconomía de la
  moneda soberana, argumentando que los gobiernos con monopolio de
  emisión no enfrentan restricciones financieras nominales, sino
  restricciones reales de recursos (inflación y capacidad instalada).
  Sostiene que la precariedad de los hogares refleja fallas en la
  distribución agregada y en la falta de redes de seguridad pública
  (como un empleo garantizado), postulando que la solución requiere
  calibración fiscal y monetaria a nivel macro.
\item
  \textbf{Visión TOC:} Sostiene que, independientemente de los arreglos
  monetarios o de la soberanía de emisor, \textbf{la restricción
  fundamental es siempre física, operativa y de capacidad productiva
  real}. Si las políticas de inyección de liquidez o los desequilibrios
  macroeconómicos expanden la masa monetaria sin elevar la capacidad
  real de producción y oferta de bienes básicos, el resultado neto es el
  incremento del costo de vida (\(OE\) macroeconómico), lo que pulveriza
  aún más el poder adquisitivo y acelera el colapso del flujo
  individual, volviendo obsoletas las intervenciones nominales.
\end{itemize}

\begin{center}\rule{0.5\linewidth}{0.5pt}\end{center}

\hypertarget{conclusiuxf3n-y-prescripciuxf3n-estratuxe9gica-5-pasos-de-enfoque-de-goldratt}{%
\subsection{5. Conclusión y Prescripción Estratégica (5 Pasos de Enfoque
de
Goldratt)}\label{conclusiuxf3n-y-prescripciuxf3n-estratuxe9gica-5-pasos-de-enfoque-de-goldratt}}

Aplicar los 5 Pasos de Enfoque de Eliyahu Goldratt al colapso de la
educación financiera exige abandonar el falso debate de la culpa
individual y rediseñar el sistema:

\begin{enumerate}
\def\labelenumi{\arabic{enumi}.}
\tightlist
\item
  \textbf{Identificar la Restricción:} Reconocer que el verdadero cuello
  de botella no es la falta de conocimientos matemáticos del ciudadano,
  sino la brecha estructural entre el costo de vida (\(OE\)) y la
  generación de ingresos reales (\(T\)).
\item
  \textbf{Explotar la Restricción:} Dejar de perder energía intentando
  exprimir centavos en presupuestos de supervivencia; aceptar que los
  márgenes locales están saturados.
\item
  \textbf{Subordinar todo lo demás a la Restricción:} Alinear las
  políticas públicas, educativas y laborales para enfocar los recursos
  en la expansión de la capacidad productiva y salarial de la base
  social, en lugar de descargar la responsabilidad en programas de
  austeridad individual.
\item
  \textbf{Elevar la Restricción:} Reestructurar los mercados, incentivar
  la innovación productiva y romper los cuellos de botella
  macroeconómicos que deprimen el valor del trabajo humano.
\item
  \textbf{Prevenir la Inercia:} Evitar que el sistema caiga de nuevo en
  el dogma de culpar al eslabón más débil (el consumidor), manteniendo
  una vigilancia constante sobre los flujos reales de valor y costo de
  la economía.
\end{enumerate}

\end{document}
